\begin{description}
\item[(12.1)] For a satellite that describes a uniform circular motion of
radius $r$ around the star and mass M, we have:
\begin{equation*}
v = \left[R_{Hodograph}\right] = \sqrt{\frac{GM}{r}} \tag*{[1 point]}
\end{equation*}

\item[(12.2)]
\begin{gather*}
\vec{a} = -\frac{Gm}{r^2}\hat{r} \tag*{[1 point]} \\
% sic: the source prints "Gm" here rather than "GM".
L = mr^2\omega \tag*{[1 point]} \\
\rightarrow\ \vec{a} = -\frac{GMm\omega}{L}\hat{r}
  = \frac{\Delta\vec{v}}{\Delta t} \\
\frac{GMm}{L}\,\frac{\Delta\theta}{\Delta t}
  = \left|\frac{\Delta v}{\Delta t}\right| \\
\Delta v = \pm\frac{GMm}{L}\Delta\theta \\
k = \frac{GMm}{L} \tag*{[2 points]}
\end{gather*}

\item[(12.3)]
\begin{gather*}
L = mv_c R \tag*{[1 point]} \\
k_c = \frac{GMm}{L} = v_c = \sqrt{\frac{GM}{R}} \tag*{[1 point]}
\end{gather*}

\item[(12.4)] In the following scheme, the hodographic circumference has
been constructed as follows:
\begin{itemize}
\item For $\theta = 0$, the velocity vector is drawn in arbitrary units (3,
  for instance) corresponding to the velocity in the periastron, $v_P$. Its
  head determines the point $A$.
\item For $\theta = \pi$, the velocity vector is drawn diametrically
  opposite in the apastron, $v_A$, also in arbitrary units. Its head
  determines point $B$. The hodograph must pass through points $A$ and $B$.
  Therefore, the radius of the hodograph must be equal to:
\begin{equation*}
R_{Hodograph} = \frac{v_P + v_A}{2} \tag*{[1 point]}
\end{equation*}
  And the ``distance'' $d$ from the star to the center of the circumference
  will be:
\begin{equation*}
d = \frac{v_P - v_A}{2} \tag*{[1 point]}
\end{equation*}
\end{itemize}
% FIGURE NEEDED: hodograph sketch for the elliptical orbit -- left: ellipse
% with the star at a focus, velocity vectors v_P (theta = 0) and v_A
% (theta = pi); right: hodograph circle on a grid through points A (head of
% v_{theta=0}) and B (head of v_{theta=pi}), star position marked off-centre
% [2 points]; 2021 T12 solution PDF, page 2.

\item[(12.5)] In the following scheme, the hodographic circumference has
been constructed as follows:
\begin{itemize}
\item For $\theta = 0$, the velocity vector is drawn in arbitrary units (4,
  for instance) corresponding to the velocity in the periastron, that is
  $v_P$. Its head determines the point $A$. Then the velocity vector is
  drawn diametrically opposite in the apastron, that is, $v_A = 0$. Its head
  determines point B that coincides with the star. The hodograph must pass
  through points A and B. Therefore, the radius of the hodograph must be
  equal to:
\begin{equation*}
R_{Hodograph} = \frac{v_P}{2} \tag*{[1 point]}
\end{equation*}
  And the distance $d$ from the star to the center of the hodographic
  circumference
\begin{equation*}
d = \frac{v_P}{2} \tag*{[1 point]}
\end{equation*}
\end{itemize}
% FIGURE NEEDED: hodograph sketch for the parabolic orbit -- left: parabola
% with the star at the focus, v_P at theta = 0 and v -> 0 as theta -> pi;
% right: hodograph circle on a grid passing through the star (point B), with
% A the head of v_{theta=0} [2 points]; 2021 T12 solution PDF, page 2.
\end{description}
